\section{The question decided, and the scope of the decision}

The sign-only method does give the requested finite-field absolute values.
Its sign is an effective multiplicity sign in a Grothendieck group of
Frobenius modules.  It is not the signature of an intersection pairing.
The missing geometric detail in the short description is supplied by
\emph{generic weak Lefschetz}, followed by cancellation through Gysin maps.
For a smooth affine $U$ of dimension $d$, the resulting effective module
contains $H_c^d(U,\mathcal F)^{\otimes k}$ and an actual Gysin kernel.
Finite Noether normalization makes its defining hypersurface compatible
with the path Fourier estimate.  This proves the middle-degree upper bound;
hyperplane induction, localization, and duality give the other degrees.

This reconstructs a proof of known weight statements, rather than claiming
that the Weil conjectures are new.  It answers the proposed identification
of methods: it uses the restriction/Gysin and evanescence machinery of
cote 16, scans 34--40, without constructing the positive involution or
establishing the cycle conjectures discussed in scans 49--50.  No new general
Hodge standard conjecture follows.  The distinction remains even in cases
where an index theorem is independently known.

The arithmetic proposed continuation fails in its unqualified form.
Norms and slopes do not choose the sign of a bilinear form; moreover,
the displayed higher-codimension formula in cote 44, scan 4, Theorem 5,
has an explicit counterexample under all its displayed hypotheses.
The valid consequence of divisibility plus the inspected slope inequality
is instead an inequality for the size of the divisibility, stated below.
For cote 34, the literal ambient group in the struck Remark 1.17 makes the
claim false; natural precise corrections are proved, including its soluble
nonabelian version.  For the final question, rational $\mathrm{SL}_2$ data
do not specify an integral lattice.  The ordinary contraction operator,
normalized Lefschetz lowering, Wilson matrices, and actual step-10 period
bounds are computed and distinguished.

\subsection*{Status convention and sources}
All propositions in Part I have complete proofs below, or are explicitly
identified as cited background.  Their check events identify author-side
checks, independent re-proofs, source inspections, and computations separately.
Part II records propositions not established here and the exact remaining
comparison data; they are not silently promoted to findings.  In particular,
an unspecified corrected Arthur/Lefschetz integral comparison is not refuted
by refuting its unnormalized version.

The stipulated Alp\"oge RH and Dirichlet-GRH statements are accepted as
theorems, exactly as the task requires.  No authenticity or build verdict is
made.  The global RH input is used for the shifted Tate factor
$\zeta(s-3)$; it is not used in the finite-field proof.  No numbered Masullo
statement is in the supplied material or used in any proof here.
The original archive \emph{images} are not supplied: the mathematical
corrections in \texttt{TRANSCRIPTION\_NOTES.tsv} are candidates for a second
reading, not claims to have deciphered handwriting.

\section{A precise sign-only theorem}
\label{sec:sign}

\subsection{Conventions and weight-free inputs}
Fix a prime power $q$, a prime $\ell\nmid q$, and an embedding
$\iota:\overline{\mathbf Q}_\ell\hookrightarrow\mathbf C$ for the
absolute values in question.  Cohomology means geometric \emph{compactly
supported} \(\ell\)-adic cohomology when the subscript $c$ is written.
$F$ denotes geometric Frobenius, so $F$ acts by $q$ on
$\overline{\mathbf Q}_\ell(-1)$.  For a rational point over
$\mathbf F_{q^n}$, its local Frobenius means Frobenius for that residue
field; if its underlying closed point has degree $a\mid n$, this is the
$(n/a)$th power of the degree-$a$ local Frobenius.

The inputs, all independent of a weight theorem, are the trace formula,
K\"unneth, affine Artin vanishing, Poincar\'e duality, and generic weak
Lefschetz.  Relevant foundational locators are SGA 4, XIV, Th\'eor\`eme 3.1
(affine vanishing), XVII, \S5.4.3 (K\"unneth), and XVIII (duality);
SGA $4\frac12$, \emph{Th\'eor\`emes de finitude}, 1.9 and 4.3 (generic
base change and duality); and the trace-formula normalization in
\cite[(1.5.1), (1.5.4), (2.3), (2.14)]{WeilI}.
The exact generic restriction theorem used is
\cite[Theorem 3.4, Corollary 3.4.1, pp.~180--181;
Appendix A.5, pp.~220--221]{KatzACT}.
Its Appendix proves the needed assertion using generic base change and
Artin vanishing; no purity assertion from the weight-theoretic parts of
that article is used.  This last separation is essential to avoiding a
circular proof.

\begin{lemma}[HS-01: effective power sums]
\label{lem:power}
Let $\lambda_1,\ldots,\lambda_s$ be distinct nonzero complex numbers and
$c_j$ positive integers.  Suppose $B>0$ and
\[
 \left|\sum_{j=1}^s c_j\lambda_j^n\right|
 \le C\exp(o(n)) B^n \qquad(n\ge1).
\]
Then $|\lambda_j|\le B$ for every $j$.
\end{lemma}
\begin{proof}
The generating power series of the sums is holomorphic on $|z|<B^{-1}$.
In a neighborhood of zero it is the rational function
\[
 \sum_j \frac{c_j\lambda_j z}{1-\lambda_j z}.
\]
At $z=\lambda_j^{-1}$ its $j$th term has a nonzero simple pole, and every
other term is regular.  Analytic continuation therefore precludes such a
pole inside the disk of convergence.  This gives the assertion.  After
coincident eigenvalues are combined, positive multiplicities prevent their
coefficients from cancelling.  The values of the power sums themselves
need not be positive.  Traces of powers depend on eigenvalues with algebraic
multiplicity, so no semisimplicity is required.
\end{proof}

\begin{theorem}[HS-02: tensor amplification]
\label{thm:tensor}
Let $V$ be a finite-dimensional Frobenius module and $a\in\mathbf R$.
For an unbounded set of positive integers $k$, suppose that an actual
Frobenius module $K_k$, a sign $e_k$, and numbers $R_{k,n}$ satisfy
\[
 e_kR_{k,n}=\operatorname{Tr}(F^n,V^{\otimes k})
                  +\operatorname{Tr}(F^n,K_k),\qquad n\ge1,
\]
and
\[
 |\iota R_{k,n}|\le C_k\exp(o_k(n))q^{n(ka+c_k)},
 \qquad c_k/k\longrightarrow0.
\]
Then every eigenvalue $\alpha$ of $F$ on $V$ satisfies
$|\iota\alpha|\le q^a$.
If the spectrum is stable under $\alpha\mapsto q^{2a}/\alpha$, it is pure
of modulus $q^a$.
\end{theorem}
\begin{proof}
For fixed $k$, the signed remainder is an effective power sum.  In
particular, $\alpha^k$ is one of its roots, with positive multiplicity.
Lemma~\ref{lem:power} gives
$|\iota\alpha|^k\le q^{ka+c_k}$.  Take $k$th roots and then let $k$ tend
to infinity.  Apply the same upper bound to $q^{2a}/\alpha$ for the final
assertion.  The limit in the extension degree $n$ is taken first, for each
fixed $k$; the constants $C_k$, dimensions of $K_k$, and later finite base
extensions may all depend on $k$.
\end{proof}

\begin{lemma}[HS-03: path Parseval estimate]
\label{lem:path}
Let $E$ be a finite field of cardinality $Q$, $A=E^d$, and $\psi$ a
nontrivial additive character.  For $k\ge2$, let $b_i:A\times A\to E$
be nondegenerate bilinear forms and let $f_i:A\to E$ be arbitrary
functions.  Put
\[
 P(x_1,\ldots,x_k)=\sum_{i=1}^{k-1}b_i(x_i,x_{i+1})
                         +\sum_{i=1}^k f_i(x_i)+c.
\]
If $|t_i(x)|\le B$, then, for every $h\ne0$,
\[
 \left|\sum_{A^k}\prod_i t_i(x_i)\psi(hP(x))\right|
       \le B^k Q^{d(k+1)/2}.
\]
The weights may depend on $h$, provided the same bound holds.  When all
$t_i=t$, character orthogonality gives
\[
 QN=M^k+R,\quad M=\sum_A t(x),\quad
 N=\sum_{P(x)=0}\prod_i t(x_i),\quad
 |R|\le B^k Q^{d(k+1)/2+1}.
\]
\end{lemma}
\begin{proof}
The matrix $K_{i,h}(x,y)=\psi(hb_i(x,y))$ satisfies
$K_{i,h}K_{i,h}^*=Q^d I$: summing the character of
$hb_i(x-x',y)$ over $y$ gives zero unless $x=x'$.  Its operator norm on
counting-measure $\ell^2(A)$ is $Q^{d/2}$.  Absorb the vertex phases into
the weights.  The path sum is a product of $k-1$ such matrices,
$k-2$ diagonal multipliers of norm at most $B$, and two endpoint vectors
of norms at most $BQ^{d/2}$.  Cauchy--Schwarz and the operator norm
inequality give $B^kQ^{d(k+1)/2}$.  A bilinear endpoint contraction instead
of a Hermitian one is handled by conjugating the first endpoint vector.
Finally
$\sum_{h\in E}\psi(hP)=Q\,1_{P=0}$; the zero mode is $M^k$, and the
$Q-1$ other modes give the claimed remainder bound.
\end{proof}

\subsection{The exact geometric sign}

\begin{proposition}[SG-08: effective Gysin remainder]
\label{prop:gysin}
Let $U/\mathbf F_q$ be smooth affine of pure dimension $d\ge1$, and
$\mathcal F$ a lisse $\overline{\mathbf Q}_\ell$-sheaf.  Put
$X=U^k$, $D=kd$, and $\mathcal F_k=\mathcal F^{\boxtimes k}$.
Suppose $i:Y\hookrightarrow X$ is a smooth affine hypersurface for which
\[
 H^a(X,\mathcal F_k^\vee)\longrightarrow
 H^a(Y,i^*\mathcal F_k^\vee)
\]
is an isomorphism for $a<D-1$ and injective for $a=D-1$.
Let $M_n$ be the trace sum of $\mathcal F$ on $U(\mathbf F_{q^n})$,
$N_{k,n}$ the trace sum of $i^*\mathcal F_k$ on $Y(\mathbf F_{q^n})$,
and $R_{k,n}=q^nN_{k,n}-M_n^k$.  Then
\begin{equation}\label{eq:effective}
 (-1)^{D-1}R_{k,n}
   =\operatorname{Tr}(F^n,H_c^d(U,\mathcal F)^{\otimes k})
      +\operatorname{Tr}(F^n,K_k),
\end{equation}
where the actual module is
\begin{equation}\label{eq:kernel}
 K_k=\ker\bigl(H_c^{D-1}(Y,i^*\mathcal F_k)(-1)
                    \longrightarrow H_c^{D+1}(X,\mathcal F_k)\bigr).
\end{equation}
\end{proposition}
\begin{proof}
Affine Artin vanishing for the dual lisse sheaf and Poincar\'e duality
give $H_c^j(U,\mathcal F)=0$ for $j<d$, and similarly the groups on
$X$ and $Y$ vanish below $D$ and $D-1$, respectively.
Dualizing the restriction in the hypothesis gives the Frobenius-equivariant
Gysin maps
\[
 H_c^j(Y,i^*\mathcal F_k)(-1)\longrightarrow
 H_c^{j+2}(X,\mathcal F_k),
\]
isomorphic for $j>D-1$ and surjective for $j=D-1$.
For clarity about the twist, the target is equivalently
$H_c^{j+2}(X,\mathcal F_k)(1)$ when the source is not twisted.

Write $T_b=H_c^{D+b}(X,\mathcal F_k)$ and
$W_a=H_c^{D-1+a}(Y,i^*\mathcal F_k)$, $a,b\ge0$.
The trace formula and $(-1)$ twist give
\[
 (-1)^{D-1}R_{k,n}
 =\sum_{a\ge0}(-1)^a\operatorname{Tr}(F^n,W_a(-1))
   +\sum_{b\ge0}(-1)^b\operatorname{Tr}(F^n,T_b).
\]
For $a\ge1$, Gysin identifies $W_a(-1)$ with $T_{a+1}$, whose sign in
the second sum is opposite.  The remaining virtual class is
$[W_0(-1)]+[T_0]-[T_1]$.  Surjectivity at the border turns this into
$[K_k]+[T_0]$.  In total K\"unneth degree $D=kd$, each of the $k$ factors
must have its minimum possible degree $d$, hence
$T_0=H_c^d(U,\mathcal F)^{\otimes k}$.
This is an equality of Frobenius-module classes, not merely dimensions.
\end{proof}

No concentration of the entire cohomology of $U$ or $Y$ in one degree was
assumed.  Higher groups cancel by actual Gysin isomorphisms; the border
surjection leaves an actual kernel.  In perverse normalization
$P=\mathcal F[d]$, form $M,N,R$ using the alternating stalk traces of
$P$ instead of those of $\mathcal F$.  The identity then has the single sign $-R$:
if $P_Y=i^*(P^{\boxtimes k})[-1]$, it is the trace of
\[
 H_c^0(U,P)^{\otimes k}\ \oplus\
 \ker\bigl(H_c^0(Y,P_Y)(-1)\longrightarrow
                         H_c^1(X,P^{\boxtimes k})\bigr).
\]
The minus sign comes from $i^*(P^{\boxtimes k})=P_Y[1]$.
The smooth-lisse proof above suffices for every conclusion below; no
additional assertion about arbitrary perverse restriction is needed.

\begin{theorem}[SG-09: the affine middle-degree bound]
\label{thm:affine}
Let $U/\mathbf F_q$ be smooth affine of pure dimension $d\ge1$, and
$\mathcal F$ lisse.  Suppose that, for a fixed $C<\infty$ and real $w$,
\[
 |\iota\operatorname{Tr}(F_{u,n},\mathcal F_{\bar u})|
       \le Cq^{nw/2}
 \quad\text{for all }n\ge1,\quad u\in U(\mathbf F_{q^n}).
\]
Every eigenvalue $\alpha$ on $H_c^d(U,\mathcal F)$ then satisfies
$|\iota\alpha|\le q^{(w+d)/2}$.
\end{theorem}
\begin{proof}
The empty case is immediate.  Choose a finite Noether normalization
$\pi:U\to\mathbf A^d$ \cite[Lemma 10.115.4, Tag 00OY]{StacksNN}.
Choose $B$ module generators of $\mathcal O(U)$ over the polynomial
coordinate ring.  Every geometric fiber has length at most $B$, hence
at most $B$ rational points over any finite extension.  Neither
\'etaleness nor constant fiber length is needed.
Choose affine coordinate generators $z_1,\ldots,z_s$ giving a closed
immersion $U\hookrightarrow\mathbf A^s$.

For each $k\ge2$, fix a nondegenerate bilinear form $b$ on
$\mathbf A^d$, and put
\[
 f(u_1,\ldots,u_k)=\sum_{i=1}^{k-1}b(\pi u_i,\pi u_{i+1}).
\]
Vary the per-vertex linear combinations $g_i=\sum_j a_{ij}z_j$ and $c$,
and consider
\begin{equation}\label{eq:pathsection}
 Y:\quad f+\sum_i g_i+c=0\quad\subset U^k.
\end{equation}
The map given by all the vertex coordinates is a closed immersion and
therefore quasi-finite.  Katz's Corollary 3.4.1, applied to this map,
the \emph{fixed} $f$, and $\mathcal F_k^\vee$, gives a dense open set of
coefficients with the restriction property of Proposition~\ref{prop:gysin}.
Generic smoothness of these hypersurfaces follows from hyperplane
Bertini for the graph closed immersion
$U^k\hookrightarrow\mathbf A^{sk+1}$, $u\mapsto(z(u),f(u))$.
The hyperplanes whose coefficient of the last coordinate is nonzero are
an open chart; rescale that coefficient to one.  Thus the smooth-good
open meets precisely the displayed family.  This argument does not
intersect a previously obtained open with an arbitrary possibly bad
fixed-coefficient slice.

Choose a point in the two good opens over some finite extension
$\mathbf F_{q^{e_k}}$.  It is permissible for $e_k$ and the coefficients
to depend on $k$.  Fix them, and write $Q=q^{e_kn}$ for any subsequent
extension degree $n\ge1$.  For $h\in\mathbf F_Q^*$ group the character
sum by $x_i=\pi u_i$, using the vertex weights
\[
 a_{i,h,n}(x)=\sum_{u\in\pi^{-1}(x)(\mathbf F_Q)}
       \iota\operatorname{Tr}(F_{u,e_kn},\mathcal F_{\bar u})
                         \psi(hg_i(u)).
\]
They have absolute value at most $BC Q^{w/2}$.  Lemma~\ref{lem:path}
applies although these weights depend on $i$ and $h$.  The zero mode
is exactly $M_{e_kn}^k$, since all vertex phases are then one.
Consequently
\begin{equation}\label{eq:parsevalbound}
 |\iota R_{k,e_kn}|\le(BC)^k
       q^{e_kn\{k(w+d)/2+d/2+1\}}.
\end{equation}
By Proposition~\ref{prop:gysin}, this remainder, up to its one sign,
is an effective power sum containing $\alpha^{e_kk}$.
Lemma~\ref{lem:power}, with $k$ fixed and all $n\ge1$, gives
\[
 |\iota\alpha|^{e_kk}
       \le q^{e_k\{k(w+d)/2+d/2+1\}}.
\]
Taking the $e_kk$th root and letting $k\to\infty$ proves the result.
No uniform bound on $e_k$ or the Betti numbers of $Y$ is required.
\end{proof}

\subsection{The weight consequences and their limits}

\begin{corollary}[SG-11: compact-support upper bounds]
\label{cor:upper}
Let $X/\mathbf F_q$ be separated of finite type, and $\mathcal F$ a
constructible $\overline{\mathbf Q}_\ell$-sheaf whose local Frobenius
eigenvalues are of $\iota$-absolute value at most $|\kappa(x)|^{w/2}$.
Then every eigenvalue on $H_c^i(X,\mathcal F)$ has $\iota$-absolute
value at most $q^{(w+i)/2}$.
\end{corollary}
\begin{proof}
First suppose $X$ is smooth affine of pure dimension $d$, with
$\mathcal F$ lisse.  In dimension zero the assertion is the local
hypothesis: induction of a residue-field Frobenius module replaces
each eigenvalue by its degree-many roots, without changing its weight.
For $d>0$, the groups in degrees $i<d$ vanish.  Theorem~\ref{thm:affine}
settles $i=d$, since the trace bound holds with the rank as $C$.
Choose another smooth generic hyperplane section $Y$; a finite base
extension is harmless.  For $i\ge d+1$ the Gysin map
\[
 H_c^{i-2}(Y,\mathcal F|_Y)(-1)\longrightarrow H_c^i(X,\mathcal F)
\]
is surjective at $i=d+1$ and isomorphic for $i>d+1$.
Induction on dimension bounds its source eigenvalues by
$q\,q^{(w+i-2)/2}=q^{(w+i)/2}$.
For a general $X$, remove nilpotents and choose a finite stratification
by smooth affine locally closed schemes on which $\mathcal F$ is lisse;
the finite field is perfect.  The localization exact sequence
\[
 H_c^i(U,\mathcal F|_U)\longrightarrow H_c^i(X,\mathcal F)
                   \longrightarrow H_c^i(Z,\mathcal F|_Z)
\]
for an open stratum union $U$ and its closed complement $Z$ shows,
by induction on the number of strata, that every middle eigenvalue
occurs in one of the adjacent groups.  Both have the desired bound.
The long exact sequence is used at degree $i$ on both sides, so there
is no extra shift loss.  A bound proved after a finite extension
descends by taking roots of the absolute values.
\end{proof}

This is the numerical weight inequality for compact support over a point,
as in \cite[Theorem 3.3.1]{WeilII}.  It is not a claim to have proved
algebraicity of arbitrary sheaf eigenvalues, all six-functor theorems,
or semisimplicity.  If the input bounds hold for every embedding under
consideration, the conclusion holds for each of them.

\begin{corollary}[SG-10: smooth proper purity and the requested Weil bounds]
\label{cor:purity}
Let $X/\mathbf F_q$ be smooth proper of pure dimension $d$, and let
$\mathcal F$ be pointwise $\iota$-pure of weight $w$ and lisse.
Then the eigenvalues on $H^i(X,\mathcal F)$ have absolute value
$q^{(w+i)/2}$.  In particular the finite-field absolute-value assertion
of Weil I holds for every smooth projective $X$, with constant
coefficients, without a Hodge standard conjecture.
\end{corollary}
\begin{proof}
Corollary~\ref{cor:upper} and properness give the upper bound.
The dual sheaf $\mathcal F^\vee$ is pointwise pure of weight $-w$.
Poincar\'e duality pairs
\[
 H^i(X,\mathcal F)\otimes
 H^{2d-i}(X,\mathcal F^\vee)\longrightarrow
                         \overline{\mathbf Q}_\ell(-d)
\]
perfectly.  Thus for every eigenvalue $\alpha$ in the first group,
$q^d/\alpha$ occurs in the second.  The latter's upper bound is
$q^{(-w+2d-i)/2}$, which gives the matching lower bound
$|\iota\alpha|\ge q^{(w+i)/2}$.
For constant coefficients the pointwise input is one, for every
embedding.  The rationality and integrality parts of Weil I are
separate established results, not deductions from the complex power-sum
lemma; the assertion reconstructed here is the requested absolute value.
\end{proof}

There is also a shorter projective-only conclusion from
Theorem~\ref{thm:affine}: choose a smooth hyperplane $H\subset X$,
use projective weak Lefschetz for $i<d$, and
$H_c^d(X\setminus H,\mathcal F)\to H^d(X,\mathcal F)\to
H^d(H,\mathcal F|_H)$ in the middle degree.  Induction gives the upper
bound, and duality gives the lower bound.  Both versions have been
checked; the compact-support proof records the stronger numerical scope
of the same mechanism.

\subsection{Why smoothness and spectral data alone are insufficient}

\begin{proposition}[HS-04 and HS-04B: the required genericity]
Smoothness of a hypersurface does not concentrate its compactly supported
cohomology in its middle degree.  Even with base middle concentration,
smoothness alone need not give the effective remainder of
Proposition~\ref{prop:gysin}.
\end{proposition}
\begin{proof}
For the first assertion take $U=\mathbf A^1$, $k=2$, constant weights,
and $Y:xy=1$.  Over $\mathbf F_{q^n}$, put $r=q^n$.
Then $M_n=r$, $N_n=r-1$, $R_n=-r$; the common top root in
$rN_n$ and $M_n^2$ cancels.  Indeed
$H_c^1(\mathbf G_m)=\overline{\mathbf Q}_\ell$ and
$H_c^2(\mathbf G_m)=\overline{\mathbf Q}_\ell(-1)$.
This example is compatible with Proposition~\ref{prop:gysin}, whose
base middle group is zero; it refutes the stronger concentration claim.

For the second assertion let $q$ be odd, let $\chi_n$ be the compatible
quadratic character of $\mathbf F_{q^n}$, extended by zero, and put
$t_n(x)=\chi_n(x^2-1)$.  This is the trace of the quadratic Kummer sheaf
on $U=\mathbf A^1\setminus\{1,-1\}$, extended by zero to the line.
The finite double cover $p:\mathbf G_m\to\mathbf A^1$,
$x=(u+u^{-1})/2$, has deck involution $u\mapsto u^{-1}$.
Its direct image splits as the constant sheaf and this extended Kummer
sheaf: at the two branch points the anti-invariant stalk is zero.
Subtracting the constant $H_c^2(\mathbf A^1)$ from the cohomology of
$\mathbf G_m$ shows that the Kummer sheaf has just
$H_c^1=\overline{\mathbf Q}_\ell$, with Frobenius eigenvalue one.
Thus $M_n=-1$ and the base is concentrated in its middle degree.

On the smooth curve $xy=1$ in $U^2$ one has
\[
 (x^2-1)(x^{-2}-1)=-\frac{(x^2-1)^2}{x^2}.
\]
Writing $\epsilon=\chi_1(-1)$ gives
\[
 N_n=\epsilon^n(q^n-3),\qquad
 R_n=(\epsilon q^2)^n-3(\epsilon q)^n-1.
\]
The expected sign $(-1)^{2-1}$ makes the coefficient of the root
$\epsilon q^2$ equal to $-1$.  Distinct finite power sums have unique
coefficients, by their generating-function poles, so $-R_n$ is not
effective.  Along this special curve the two ramification factors
cancel, leaving a nonzero top compact cohomology group.  The generic
restriction hypothesis fails.  Generic weak Lefschetz in the actual
Fourier-compatible family, rather than smoothness by itself, excludes
precisely this defect.
\end{proof}

\begin{proposition}[HS-05: weights and parity do not specify a polarization]
The Frobenius spectrum, its cohomological parity, and a Frobenius-compatible
nondegenerate symmetric form do not by themselves imply that form is
positive definite.
\end{proposition}
\begin{proof}
Place $V=\mathbf Q^2$ in degree $2r$ and let $F=q^r I$.
The forms with matrices $\operatorname{diag}(1,1)$ and
$\operatorname{diag}(1,-1)$ have the same compatibility
$b(Fx,Fy)=q^{2r}b(x,y)$ and identical spectral data, but signatures
$(2,0)$ and $(1,1)$.  This is a countermodel to the formal inference,
not a geometric counterexample to the Hodge standard conjecture.
Likewise a nontrivial Jordan block and its scalar semisimplification
have the same traces of all powers; the power-sum proof contains no
semisimplicity conclusion.
\end{proof}

\subsection{What was located in the supplied and retrieved Lean sources}
\label{sec:lean}
\paragraph{HS-06 (source-level conclusion).}
The short prose description is in Alp\"oge's \emph{Remark}, p.~3,
footnote 3, referring to timeline steps 3 and 5.  The supplied
\texttt{ALPOGE\_Main.lean} and \texttt{ALPOGE\_Statements.lean} are
wrappers for the stipulated global RH/Dirichlet-GRH statements.
For the method comparison we subsequently obtained the public release
archive specified in \texttt{evidence/LEAN\_MIRROR\_PROVENANCE.txt}; both
wrappers are byte-identical to their supplied versions.

The relevant retrieved file
\path{Development/EtaleCohomology/WeilConjectures/FermatCount.lean}
defines, in namespace \texttt{LateralTheory.EWeil}, the propositions
\texttt{Base}, \texttt{Goal}, \texttt{WeakAt}, \texttt{Tensor},
\texttt{FermatCount}, \texttt{HypTrace}, and others.
\texttt{Goal} has an explicit \texttt{Base} premise and concludes
$\|\iota\alpha\|^2=q^n$ for smooth projective cohomology;
\texttt{WeakAt} allows the interval $[n-3,n+3]$, and
\texttt{Tensor} sharpens such a dimension-independent loss.
The occurrence of a proposition definition or an implication between
propositions is not itself a closed theorem supplying all premises.
A targeted textual search of all 175,187 source files, and closer
declaration-level inspection of the 24 files in this directory and
the named divisibility/slope directories, are documented separately.
This is not a Lean kernel certificate or a proof of absence of a differently
named theorem.  The mathematical proof in this paper stands on its
explicit geometric inputs; no unlocated formal theorem is used to fill
a gap.  The arithmetic section gives exact declaration locators for the
steps 10 and 13 statements actually inspected.
